\bookchapter{Chapter 17: Interconnection Between Phone and Other Devices}

\chapterintro{Android + Windows (Link to Windows)}

\booksection{: Microsoft has an app called “Link to Windows”, built into Samsung phones, and other Android phones can download it from the app store. After pairing, you can directly view phone notifications, reply to texts, browse photos, and even make and receive calls on your computer. I've used it for two years—during meetings I don't have to keep looking down at my phone; a notification popping up on the computer screen is enough.}

Brand multi-screen Collaboration

\begin{itemize}
  \item \textbf{: Huawei's "Multi-Screen Collaboration" and Xiaomi's "Xiaomi Smart Connect" go even further—directly project your phone screen onto your computer. You can use the mouse to drag files from phone to computer, and also browse TikTok and reply to WeChat on the computer. Cross-device copy-paste and file dragging all work, with a very smooth experience.}2. Screen Casting to TV: Watch Together with the Whole Family
  \item \textbf{When I went home for Chinese New Year, my cousin wanted to play her graduation trip video from her phone on the TV. After fumbling for a long time, we found it was actually very simple:}iPhone TV
  \item \textbf{: Swipe down from the top right to open Control Center, tap the icon with two overlapping squares (Screen Mirroring), find your TV's name, and tap to connect. The prerequisite is that the TV and phone are on the same Wi-Fi, and the TV supports AirPlay (most smart TVs do).}Android TV
\end{itemize}

\booksection{: Each brand calls it differently—Huawei calls it “Wireless Projection”, Xiaomi calls it “Cast”, OPPO/vivo calls it “Multi-Screen Interaction”. Usually findable in the quick settings panel; if not, search “screen cast” in Settings.}

Direct cast from video apps

\begin{itemize}
  \item \textbf{: If you just want to watch iQiyi, Bilibili, or YouTube on TV, it's even simpler—find a small rectangle icon with a Wi-Fi symbol on the playback interface, tap it, select the TV, and the video goes over. At this point, the phone can do other things while the video continues playing on the TV, without interfering with each other.}What if screen casting fails
  \item \textbf{: 90\% of the time, it's because they're not on the same Wi-Fi. The other 10\% is that the TV's screen casting function isn't enabled—look for the "Multi-Screen Interaction" or "AirPlay" switch in TV settings.}3. With Smart Watches/Bands: A Health Manager on Your Wrist
  \item \textbf{I bought my mom a fitness band. At first she said "I don't exercise, why buy this?" After two weeks of use, she discovered she was only sleeping five hours a day, and started proactively adjusting her schedule. This is the most practical value of watches/bands—turning invisible health data into visible numbers.}Notification filter
  \item \textbf{: When your phone keeps vibrating in your bag or pocket, just raise your wrist to take a look—is it a WeChat from the boss or a Pinduoduo promotion? Judge in a second, and let messages not worth pulling out the phone for go.}Health data sync
\end{itemize}

\booksection{: The watch monitors heart rate, blood oxygen, sleep, and steps 24 hours a day, and this data is automatically aggregated into the health app on your phone. Over time, you can see your trends—have you been walking too little this month, has your sleep quality improved this week.}

Transit card/access card sync

\begin{itemize}
  \item \textbf{: Transit cards and access cards bound on your phone will mostly auto-sync to your watch. When taking the subway, you don't need to pull out your phone—just tap your wrist against the turnstile, and you really get a "the future is here" feeling.}Find phone lifesaver
  \item \textbf{: Watches generally have a "find phone" feature—tap it and the phone will ring. I use this feature far more than expected—I just put the phone down somewhere and forget where.}4. With Smart Home: The Phone Is the Home's "Master Switch"
  \item \textbf{Smart home sounds mysterious, but it boils down to three things: remote control via phone, set automation rules, voice control.}Let me talk about a few of my own actual use cases, so you can see how practical this thing is:
  \item \textbf{Air conditioner}: On the way home in summer, use your phone to turn on the AC in advance—when you walk in, it's already cool. Buy an AC companion socket for a hundred-something yuan to achieve this, no need to replace the AC.
\end{itemize}

\booksection{Lights}

: Lying in bed before sleep, too lazy to get up to turn off the lights—just tap on your phone or say “Hey Siri/Xiaoai, turn off the lights”.

Robot vacuum

\begin{itemize}
  \item \textbf{: Before going out, tap "Start Cleaning" on your phone—when you come back, the floor is clean. You can also set a daily scheduled cleaning.}Camera
  \item \textbf{: If you have elderly, children, or pets at home, you can check on the home situation anytime from your phone.}Which ecosystem to choose? One-sentence summary:
  \item \textbf{Using iPhone prioritize devices that support Apple HomeKit (a bit more expensive but better privacy protection).}Many devices, want cost-effectiveness Xiaomi Mi Jia ecosystem is the most complete—lights, sockets, robot vacuums, everything.
  \item \textbf{Using Huawei phone prioritize HarmonyOS Connect, best experience with Huawei devices.}5. Bluetooth Headphones and Speakers: Wireless Freedom
\end{itemize}

Switching from wired to Bluetooth headphones was, for me, one of the most noticeable upgrades in phone experience. No more pulling out a tangled mess from the bag, no more worrying about the cable catching on door handles.

\begin{itemize}
  \item Pairing is simple
  \item : Headphones generally enter pairing mode automatically when first turned on (flashing light). Turn on Bluetooth on your phone to find them, tap to connect. After that, they auto-connect every time you take them out—nothing to do.
  \item Same brand has bonuses
\end{itemize}

\booksection{: AirPods with iPhone will auto-pop up to show battery, and can seamlessly switch between multiple Apple devices. Huawei FreeBuds with Huawei phones, Xiaomi earbuds with Xiaomi phones, also have similar pop-ups and optimizations.}

Bluetooth speakers

\begin{itemize}
  \item \textbf{: For family gatherings or outdoor camping, a Bluetooth speaker can double the atmosphere. After connecting, your phone becomes your DJ booth.}Watch out for latency
  \item \textbf{: When playing games with Bluetooth headphones, there may be slight latency (you hear the gunshot half a second after it fires). If you're sensitive to latency, either use wired headphones or buy Bluetooth headphones that support "low latency mode".}6. With Car Systems: A Safe Partner While Driving
  \item \textbf{Looking at your phone while driving is illegal and life-threatening. But sometimes you really need navigation, phone calls, and music. This is where CarPlay (Apple) and Android Auto (Android) come in handy.}How to connect
  \item \textbf{: If your car supports it (most new cars in recent years do), use a data cable to plug your phone into the car's USB port, and the car screen will automatically switch to the CarPlay or Android Auto interface. Some new cars support wireless connection, even more convenient.}What it can do
\end{itemize}

\begin{tcbinfo}{UWB (Ultra-Wideband) Technology — More Precise "Spatial Awareness" Than Bluetooth}
In addition to Bluetooth, new phones also contain a UWB (Ultra-Wideband) chip. Its capability is\textbf{precise distance measurement and spatial positioning}— it can calculate how far another device is and in which direction, with centimeter-level precision. Bluetooth can only tell you “it's nearby,” while UWB can tell you “it's 1.2 meters to your front-left.” Common applications:

\begin{itemize}
  \item \textbf{AirTag / Precise Finding}: When an iPhone finds an AirTag, it shows directional arrow navigation like "7 meters away, go left" — this relies on UWB. Regular Bluetooth can only tell you whether it's within range.
  \item \textbf{Digital Car Keys}: iPhones and Android flagships with UWB can serve as car keys — automatically unlock when approaching the vehicle, lock when walking away, and distinguish between driver and passenger positions to prevent relay attacks.
  \item \textbf{Indoor Navigation}: "Aisle-level" navigation in malls and airports via UWB beacons, much more accurate than GPS.
  \item \textbf{Supported Models}: iPhone 11 and later all have UWB chips (U1/U2); on the Android side, some flagships from Xiaomi, Samsung, vivo, and OPPO are also equipped.
\end{itemize}

\end{tcbinfo}

\begin{tcbinfo}{LE Audio and Auracast — The New Generation Bluetooth Audio Standard}
The classic Bluetooth audio (A2DP) used for twenty years is being replaced by\textbf{LE Audio}(Low Energy Audio), the new generation audio standard from Bluetooth 5.2 onwards:

\begin{itemize}
  \item \textbf{Lower Power Consumption, Better Sound Quality}: Uses the all-new LC3 codec, with better sound quality than the old SBC standard and lower power consumption — earbud battery life is noticeably extended.
  \item \textbf{Dual-Path Audio}: Left and right earbuds transmit independently — switching to single-ear mode won't cause disconnection.
  \item \textbf{Auracast Broadcast Audio}: The signature feature of LE Audio. A single phone/TV can simultaneously broadcast audio to\textbf{countless pairs of}earbuds — airport waiting areas, museums, gyms, and classrooms can all use a single transmitter for multiple people to listen with their own earbuds.
  \item \textbf{Current status}: Flagship earbuds and phones released in 2024-2026 generally support LE Audio, but older devices require hardware upgrades and can't be software-unlocked. When buying new earbuds, check for the LE Audio/Auracast label.
\end{itemize}

\end{tcbinfo}

\booksection{: Navigation (much better than built-in car navigation), phone calls (answer from steering wheel), music (music apps on your phone directly cast to the car), voice WeChat (Siri or Xiaoai helps read messages and reply).}

Old cars can use it too

\begin{itemize}
  \item \textbf{: If your car is older and doesn't support it, buy an external CarPlay screen for a few hundred yuan, stick it on the dashboard—the experience is about the same.}7. Security and Privacy in Integration
  \item \textbf{The more devices connected, the larger the surface area for privacy leaks. A few points to note:}Before screen casting in public, first check whether your phone notification bar has any private messages—don't let the whole family see what they shouldn't.
  \item \textbf{Before giving or selling a device to someone else, unbind it first—go to phone settings, find that device, and select "Remove" or "Unpair".}For smart home cameras, buy branded, encrypted ones—don't buy no-name ones for a few dozen yuan. There's been too much news about cameras being hacked for live streaming.
\end{itemize}

\booksection{Try It: Cast Your Phone Screen to TV and Set Up a Smart Scene}

Make sure your phone and TV are on the same Wi-Fi.

\begin{itemize}
  \item Find "Screen Mirroring" or "AirPlay" on the TV and turn it on.
  \item On your phone: iPhone taps "Screen Mirroring" from Control Center; Android swipes down to find "Wireless Projection".
  \item Cast your favorite photo to the TV for your family to see.
\end{itemize}

\begin{tcbpractice}{If you have smart home devices, try creating a "Leave Home Mode" in the phone app: one-tap to turn off lights, turn off AC, and start the robot vacuum.}
\begin{enumerate}
  \item Chapter 18: Data Backup and Migration
  \item A neighbor of mine dropped her phone into a pool, and by the time she fished it out, it wouldn't turn on. What hurt her most wasn't the phone itself, but the three years of unbacked-up photos—every record of her child from birth to learning to walk. After hearing her story, I felt terrible, because this was completely avoidable. In this chapter, we'll seriously put insurance on your data.
  \item 1. First, Figure Out: What in Your Phone Is Gone Forever If Lost
  \item Different data has different "survival states"—you need to distinguish them:
  \item Stored locally on the phone, truly gone if lost
\end{enumerate}

\end{tcbpractice}

\begin{tcbinfo}{Case 1: Battery Life Optimization in Practice}
\begin{itemize}
  \item \textbf{Scenario}: Xiaomi 14 user with poor battery life — out of power by midday.
  \item \textbf{Steps}: (1) Settings Battery Check power consumption ranking (2) Identify abnormal background apps (3) Switch from 5G to 4G (4) Restrict background activity (5) Enable power saving mode (6) Remove case while charging for heat dissipation.
  \item \textbf{Expected Result}: Battery life extended by 30-50\%, battery health data readable.
\end{itemize}

\end{tcbinfo}
